% Generated draft: requires mathematical and source review.
\documentclass[11pt]{article}
\usepackage[margin=25mm]{geometry}
\usepackage{fontspec}
\setmainfont{Latin Modern Roman}
\usepackage{amsmath,amssymb,mathtools,unicode-math}
\setmathfont{Latin Modern Math}
\usepackage{enumitem,needspace,xurl,hyperref}
\setlength{\parindent}{0pt}
\setlength{\parskip}{5pt}
\setlength{\emergencystretch}{3em}
\newcommand{\N}{\mathbb{N}}
\newcommand{\Z}{\mathbb{Z}}
\newcommand{\Q}{\mathbb{Q}}
\newcommand{\R}{\mathbb{R}}
\newcommand{\C}{\mathbb{C}}
\newcommand{\F}{\mathbb{F}}
\newcommand{\A}{\mathbb{A}}
\newcommand{\PP}{\mathbb{P}}
\newcommand{\E}{\mathbb{E}}
\newcommand{\eps}{\varepsilon}
\newcommand{\dd}{\,\mathrm{d}}
\begin{document}
\section*{UnsolvedMath ID: 701066}
\textbf{Source problem code:} catalogue entry
\subsection{Definitions and mathematical statement}
In the usual setting of formations of finite groups, a \emph{formation} is a nonempty class \(\mathcal F\) of finite groups, closed under isomorphisms, homomorphic images, and finite subdirect products. Here a subdirect product of \(G_1,\ldots,G_r\), with \(r\geq 1\), is a subgroup \(D\leq G_1\times\cdots\times G_r\) whose projection onto every factor is surjective. For a finite group \(G\), define its generated formation by
\[
\operatorname{form}(G)=\bigcap\{\mathcal F: \mathcal F\text{ is a formation of finite groups and }G\in\mathcal F\}.
\]
A \emph{subformation} of \(\operatorname{form}(G)\) is a formation \(\mathcal H\) with \(\mathcal H\subseteq\operatorname{form}(G)\).

The supplied assertion is that, for every finite group \(G\), the formation \(\operatorname{form}(G)\) has only finitely many \(S_n\)-closed subformations. The source record does not define \(S_n\)-closed or explicitly specify the ambient universe for formations, so a completely unqualified formalization is unavailable from the supplied record alone.

If \(S_n\)-closed means closed under taking subnormal subgroups, the assertion has the following precise formulation. Write \(L\mathrel{\trianglelefteq\trianglelefteq}K\) when there are an integer \(r\geq 0\) and a chain
\[
L=L_0\trianglelefteq L_1\trianglelefteq\cdots\trianglelefteq L_r=K.
\]
Call a formation \(\mathcal H\) subnormal-subgroup-closed when
\[
\forall K\in\mathcal H\;\forall L\leq K,\qquad
L\mathrel{\trianglelefteq\trianglelefteq}K\ \Longrightarrow\ L\in\mathcal H.
\]
Under this interpretation, the proposed statement is
\[
\forall\text{ finite groups }G,\qquad
\left|\left\{\mathcal H\subseteq\operatorname{form}(G):
\begin{array}{l}
\mathcal H\text{ is a formation of finite groups},\\
\mathcal H\text{ is subnormal-subgroup-closed}
\end{array}
\right\}\right|<\infty.
\]
This explicit formulation is conditional on the interpretation of \(S_n\) and the finite-group setting. In particular, it does not treat \(n\) as a separately quantified integer or impose a bound on subnormal defect; the supplied record does not settle those notational questions.
\subsection{Short English statement}
The proposed assertion says that the formation generated by each finite group contains only finitely many \(S_n\)-closed subformations. If \(S_n\)-closure means closure under subnormal subgroups, the displayed formulation states this explicitly. The meaning of \(S_n\) and the ambient setting require confirmation from the source.
\subsection{Sources}
\begin{itemize}
\item Supplied local UnsolvedMath registry record, permanent ID \(701066\), exact source problem code \texttt{catalogue entry}, titled ``Kourovka Notebook Problem 14.81.'' The supplied registry labels the problem \emph{open}; this is a source label, not an independently verified mathematical or literature assessment.
\item Attributed source: \emph{The Kourovka Notebook (Unsolved Problems in Group Theory)}, 14th issue (1999), Problem 14.81. Proposed by A. N. Skiba and L. A. Shemetkov.
\item The supplied background also identifies the 21st issue edition (2026), arXiv:1401.0300v48. The collection description credits E. I. Khukhro and V. D. Mazurov as editors. No source text defining \(S_n\), external URL, or literature-verification date is supplied; the cited work has not been independently consulted for this draft.
\end{itemize}
\end{document}
